% ID: SEL-D0066-OB-E01
% ESTADO: validada. S066.
\textbf{a)} Las cotas son $8-4x\le y\le(4x+16)/3$ y $x\le2$. Las dos cotas de $y$ son compatibles cuando $x\ge1/2$. Las rectas oblicuas se cortan en $(1/2,6)$. En $x=2$ los extremos son $(2,0)$ y $(2,8)$. El recinto es el triángulo cerrado de esos vértices.
\begin{center}
\setlength{\unitlength}{0.9mm}
\begin{picture}(90,60)(-8,-8)
\qbezier(0,0)(35,0)(70,0)
\qbezier(0,0)(0,24)(0,48)
\qbezier(15,30)(37.5,15)(60,0)
\qbezier(60,0)(60,20)(60,40)
\qbezier(60,40)(37.5,35)(15,30)
\put(72,0){$x$}\put(2,48){$y$}
\put(2,34){\scriptsize $(1/2,6)$}
\put(62,39){\scriptsize $(2,8)$}
\put(57,-6){\scriptsize $(2,0)$}
\end{picture}
\end{center}
\textbf{b)} En $(1/2,6),(2,0),(2,8)$, los valores son $-9/2,6,-2$. Por linealidad, el mínimo es $\boxed{-9/2}$ sólo en $(1/2,6)$ y el máximo es $\boxed{6}$ sólo en $(2,0)$.
